\chapter{Introducción}
\label{chap:intro}

\section{Motivación del proyecto}

Muchos mecanismos de elevación no presentan una carga constante al actuador. En una barra
articulada que se levanta tirando de un cable, la tensión necesaria para sostener la carga depende
del ángulo que forman el cable y la barra; al girar la barra, esa geometría cambia y con ella la
tensión. Si el cable se enrolla en un tambor cilíndrico, el par que el motor debe aportar, producto
de la tensión por el radio, sigue fielmente esas variaciones, y el actuador hay que dimensionarlo
para el peor punto del recorrido aunque la mayor parte del tiempo trabaje muy por debajo.

Una solución clásica es cambiar la relación de transmisión a lo largo del recorrido haciendo que el
radio del tambor no sea constante. Es la idea del \emph{fusee} de la relojería del siglo XV, un cono
con una ranura helicoidal que compensaba la pérdida de par del muelle real a medida que se
desenrollaba \cite{landes1983}. La misma idea reaparece hoy en los tambores de radio variable de la
robótica de cables \cite{seriani2016,bieber2023} y en los sistemas pasivos de compensación de
gravedad \cite{shen2022,endo2010}. En este trabajo, al tambor en espiral se le llama \emph{caracola}
por su forma.

El mecanismo objeto del proyecto es precisamente una caracola que, accionada por un servomotor paso
a paso de lazo cerrado, enrolla un cable y levanta una barra articulada con una masa de
\SI{2}{\kilo\gram}. La hipótesis de partida es que el perfil de radio variable aplana el par que ve
el motor frente a la variación de la tensión del cable. Comprobarlo exige dos cosas: un modelo que
diga qué cabe esperar y un sistema de medida capaz de registrar a la vez la posición del eje, la
tensión del cable y el par del motor en ensayos repetibles. Construir ese sistema, y usarlo para
contrastar el modelo, es el contenido de este Trabajo Fin de Grado.

\section{Objetivos}

El objetivo general es diseñar, construir y validar un banco de ensayos para el mecanismo de
caracola: la electrónica de accionamiento, control y adquisición, el software de supervisión y el
modelo con el que se interpretan las medidas.

De él se derivan los siguientes objetivos específicos:

\begin{enumerate}[label=\textbf{O\arabic*.}]
  \item Desarrollar un modelo estático del mecanismo que relacione la geometría, la tensión del
        cable y el par en el tambor, y deducir de él el perfil de radio que hace constante el par.
  \item Diseñar una placa de circuito impreso que adapte el microcontrolador al entorno de la
        máquina: bus RS-485, entradas digitales protegidas, entradas analógicas acondicionadas,
        salidas de potencia y conexión de células de carga.
  \item Desarrollar el firmware del microcontrolador, capaz de gobernar el servoaccionamiento por
        Modbus RTU, leer la célula de carga, ejecutar el \emph{homing} y los ciclos de ensayo, y
        vigilar los límites de seguridad con una temporización determinista.
  \item Definir un protocolo de comunicación robusto entre el microcontrolador y el ordenador, que
        transporte la telemetría por lotes y las órdenes con detección de errores.
  \item Desarrollar una aplicación de supervisión web para operar la máquina, visualizar los datos
        en tiempo real, grabar cada ensayo en un fichero reproducible, configurar la geometría del
        mecanismo y comparar cada ensayo con el modelo.
  \item Caracterizar el banco (temporización, repetibilidad y comportamiento de la medida) y
        analizar los ensayos para contrastar el modelo y cuantificar el efecto de la caracola.
  \item Documentar los problemas técnicos encontrados y las soluciones adoptadas, así como la
        planificación, el presupuesto y el impacto del trabajo.
\end{enumerate}

\section{Materiales utilizados}

\begin{itemize}
  \item \textbf{Mecanismo}: caracola de radio variable, cable, barra articulada con una masa de
        \SI{2}{\kilo\gram} y dos finales de carrera
        \completar{indicar si la caracola y la estructura se han diseñado y fabricado en este TFG o
        proceden de un trabajo previo}.
  \item \textbf{Accionamiento}: servomotor paso a paso de lazo cerrado con accionamiento Modbus RTU
        sobre RS-485 \completar{marca y modelo}.
  \item \textbf{Electrónica}: placa adaptadora propia diseñada en KiCad \cite{kicad} alrededor de un
        módulo Raspberry Pi Pico~2 \cite{rp2350}, con transceptor THVD1400 \cite{thvd1400},
        interruptor de potencia TPS4H000-Q1 \cite{tps4h000} y amplificadores MCP6294
        \cite{mcp629x}.
  \item \textbf{Medida de fuerza}: célula de carga con acondicionador ZSC31014 \cite{zsc31014}
        montado en una placa \emph{Load Cell 2 Click} \cite{mikroelc2}.
  \item \textbf{Software}: entorno Arduino con el núcleo Arduino-Pico \cite{arduinopico} y la
        biblioteca ModbusMaster \cite{modbusmaster} para el firmware; Python con FastAPI
        \cite{fastapi}, pySerial \cite{pyserial} y Chart.js \cite{chartjs} para la supervisión;
        matplotlib \cite{matplotlib} para el análisis, y git para el control de versiones.
\end{itemize}

\section{Alcance y limitaciones}

Entra en el alcance todo el sistema electrónico y de software (placa, firmware, protocolo,
supervisión y análisis), el modelo del mecanismo y la campaña de ensayos. El diseño mecánico se
toma como dado.

Las limitaciones conocidas, que se discuten en los capítulos correspondientes, son:

\begin{itemize}
  \item La posición del eje no se lee del codificador del accionamiento, sino que se obtiene
        integrando la velocidad que devuelve el propio servo (sección~\ref{sec:fw-posicion}).
  \item La célula de carga no está calibrada con masas patrón; los ensayos se guardan en cuentas
        brutas y la conversión a newtons es provisional (sección~\ref{sec:res-limitaciones}).
  \item El par se obtiene en las unidades internas del accionamiento, cuya escala física no se ha
        confirmado, de modo que los resultados se expresan en variaciones relativas.
  \item Las cotas del mecanismo no se han medido: los parámetros geométricos proceden de un ajuste
        sobre los propios ensayos (capítulo~\ref{chap:resultados}).
  \item Los finales de carrera y el límite de fuerza se implementan por software; no hay una cadena
        de seguridad cableada (sección~\ref{sec:arq-seguridad}).
\end{itemize}

\section{Metodología}

El desarrollo ha sido incremental y guiado por pruebas. Cada subsistema (bus RS-485, célula de
carga, finales de carrera, entradas y salidas de la placa) se validó primero con un programa de
prueba aislado, que se conserva en el repositorio (anexo~\ref{anx:repo}), antes de integrarlo en el
firmware principal. El código se ha versionado con git desde el primer día, y ese historial ha
servido además para reconstruir la planificación real (capítulo~\ref{chap:gestion}). Los ensayos se
guardan en crudo, con las constantes de conversión en la cabecera de cada fichero, de modo que
cualquier resultado puede recalcularse sin repetir el ensayo; las figuras de resultados de esta
memoria se generan con programas reproducibles a partir de esos ficheros.

\section{Estructura del documento}

\begin{itemize}
  \item El capítulo~\ref{chap:marco} resume los conceptos necesarios: accionamientos, medida de
        fuerza con galgas, buses de campo y microcontroladores.
  \item El capítulo~\ref{chap:estado} sitúa el trabajo frente a los tambores de radio variable y los
        mecanismos de compensación publicados.
  \item El capítulo~\ref{chap:modelo} desarrolla el modelo del mecanismo y el método experimental
        con el que se contrasta.
  \item El capítulo~\ref{chap:arquitectura} recoge los requisitos y la arquitectura del banco.
  \item Los capítulos~\ref{chap:hardware} a~\ref{chap:software} describen el hardware, el firmware,
        las comunicaciones y el software de supervisión.
  \item El capítulo~\ref{chap:problemas} reúne los problemas técnicos y sus soluciones.
  \item El capítulo~\ref{chap:resultados} presenta la caracterización del banco y los resultados de
        las dos campañas de ensayos.
  \item El capítulo~\ref{chap:gestion} trata el ciclo de vida, la planificación y el presupuesto, y
        el capítulo~\ref{chap:conclusiones} las conclusiones y el trabajo futuro.
\end{itemize}

Los anexos contienen la asignación de pines, el protocolo, los registros Modbus, el formato de los
ficheros de ensayo, el manual de uso y la estructura del repositorio.
